\section*{第 \ref{ch:b1:extent-q-and-k} 章——描述化学体系：反应进度、活度、Q 与 K}

\begin{solution}{exo:b1:extent-q-and-k:1}
强度量：温度、物质的量浓度、密度、压强、摩尔
分数。广延量：体积、物质的量、质量。
\end{solution}

\begin{solution}{exo:b1:extent-q-and-k:2}
$p(\ce{N2}) = 0.7808 \times 1.013 = \qty{0.791}{bar}$，$p(\ce{O2}) =
0.2095 \times 1.013 = \qty{0.212}{bar}$，$p(\ce{Ar}) = 0.0093 \times 1.013
= \qty{0.0094}{bar}$。
\end{solution}

\begin{solution}{exo:b1:extent-q-and-k:3}
$n(\ce{H2}) = 3.0 - 2\xi$，$n(\ce{O2}) = 2.0 - \xi$，$n(\ce{H2O}) =
2\xi$。$\xi_{\max} = \min(3.0/2,\ 2.0/1) = \qty{1.5}{mol}$：氢气为限量反应物。
最终状态：\ce{H2} 0，\ce{O2} \qty{0.5}{mol}，\ce{H2O}
\qty{3.0}{mol}。
\end{solution}

\begin{solution}{exo:b1:extent-q-and-k:4}
(a) $Q = p_{\ce{CO2}}/p^\circ$。(b) $Q = (c^\circ)^2/([\ce{Ag+}][\ce{Cl-}])$。
(c) $Q = (p_{\ce{NH3}}/p^\circ)^2/[(p_{\ce{N2}}/p^\circ)(p_{\ce{H2}}/p^\circ)^3]$。
(d) $Q = [\ce{NH4+}][\ce{OH-}]/([\ce{NH3}]\,c^\circ)$，水的活度为 1。
\end{solution}

\begin{solution}{exo:b1:extent-q-and-k:5}
$\mathrm{A \rightleftharpoons C}$ 是上两式之和：$K = 10 \times 0.5 = 5$。$\mathrm{C \rightleftharpoons A}$：$1/5 =
0.2$。$\mathrm{2\,A \rightleftharpoons 2\,C}$：$5^2 = 25$。
\end{solution}

\begin{solution}{exo:b1:extent-q-and-k:6}
$Q = 5.0 \times 10^{-5} < K^\circ$：正向进行。$Q = 0.10 > K^\circ$：逆向进行。
仅有反应物：$Q = 0 < K^\circ$，正向进行。
\end{solution}

\begin{solution}{exo:b1:extent-q-and-k:7}
$n(\ce{A2}) = 1 - \xi$，$n(\ce{A}) = 2\xi$，总量为 $1 + \xi$。取 $p =
p^\circ$，$Q = \dfrac{(2\xi/(1+\xi))^2}{(1-\xi)/(1+\xi)} = \dfrac{4\xi^2}{1 -
\xi^2}$。由 $Q = 0.25$ 得 $4.25\,\xi^2 = 0.25$，$\xi = \qty{0.243}{mol}$。
\end{solution}

\begin{solution}{exo:b1:extent-q-and-k:8}
设 $x = [\ce{C}]$：$x/(0.10 - x)^2 = 100$，故 $100x^2 - 21x + 1 = 0$ 且
$x = (21 - \sqrt{41})/200 = \qty{0.0730}{mol/L}$（另一根超过
0.10）。$[\ce{A}] = [\ce{B}] = \qty{0.0270}{mol/L}$，$[\ce{C}] =
\qty{0.0730}{mol/L}$，$\tau = 0.73$。
\end{solution}

\begin{solution}{exo:b1:extent-q-and-k:9}
由各物质的 $n$ 可得：$Q = \xi^2/(n - \xi)^2 = K^\circ$，故 $\xi/(n - \xi) =
\sqrt{K^\circ}$，且 $\tau = \xi/n = \sqrt{K^\circ}/(1 + \sqrt{K^\circ})$。
$\tau = 0.99$ 要求 $\sqrt{K^\circ} = 99$，$K^\circ \approx \num{9.8e3}$：
约 $10^4$，即定量反应通常的阈值。
\end{solution}

\begin{solution}{exo:b1:extent-q-and-k:10}
平衡时 $p_{\ce{CO2}} = 0.20\,p^\circ$，即 $n(\ce{CO2}) =
pV/RT = 0.20 \times 10^5 \times 0.0100/(8.314 \times 1100) =
\qty{0.0219}{mol}$。若碳酸盐为 \qty{0.010}{mol}，则无法
达到此值：全部分解，$p_{\ce{CO2}} = 0.010 \times 8.314 \times
1100/0.0100 = \qty{9.15e3}{Pa} = \qty{0.091}{bar}$，$Q < K^\circ$，且
最终状态（\qty{0.010}{mol} \ce{CaO}，\qty{0.010}{mol} \ce{CO2}，无
\ce{CaCO3}）并非平衡态。若为 \qty{0.050}{mol}，则达平衡：
\qty{0.0219}{mol} \ce{CO2} 与 \ce{CaO}，剩余 \qty{0.0281}{mol} \ce{CaCO3}。
\end{solution}

\begin{solution}{exo:b1:extent-q-and-k:11}
平衡时 $n(\ce{B}) = 2\,n(\ce{A})$ 且 $n(\ce{C}) = 3\,n(\ce{A})$
（体积相互抵消）。于是 $6\,n(\ce{A}) = 1.0$：$n(\ce{A}) =
\qty{0.167}{mol}$，$n(\ce{B}) = \qty{0.333}{mol}$，$n(\ce{C}) =
\qty{0.500}{mol}$。
\end{solution}

\begin{solution}{exo:b1:extent-q-and-k:12}
等物质的量：$\xi^2/(0.50 - \xi)^2 = 4$，$\xi/(0.50 - \xi) = 2$，
$\xi = \qty{0.333}{mol}$，产率 $0.333/0.50 = 67\,\%$。若醇为 \qty{1.0}{mol}
时：$\xi^2 = 4(0.50 - \xi)(1.0 - \xi)$，$3\xi^2 - 6\xi + 2 = 0$，
$\xi = (6 - \sqrt{12})/6 = \qty{0.423}{mol}$，产率 85\,\%：过量的一种反应物
会提高另一反应物所生成产物的产率。
\end{solution}

\begin{solution}{pb:b1:extent-q-and-k:1}
\textbf{1.} $x(\ce{CH4}) = 1.00/4.00 = 0.250$；$x(\ce{H2O}) = 0.750$。
\textbf{2.} \qty{7.5}{bar} 和 \qty{22.5}{bar}。
\textbf{3.} 质量为 16.04 和 \qty{54.06}{g}：$w(\ce{CH4}) = 0.229$，
$w(\ce{H2O}) = 0.771$。
\textbf{4.} 摩尔分数、质量分数、分压（以及
总压）；物质的量均为广延量。
\textbf{5.} $a = p_i/p^\circ$：甲烷为 7.5，水蒸气为 22.5。
\textbf{6.} \ce{CH4} $1.00 - \xi$，\ce{H2O} $3.00 - \xi$，\ce{CO} $\xi$，
\ce{H2} $3\xi$，总量 $4.00 + 2\xi$。
\textbf{7.} $\xi_{\max} = \qty{1.00}{mol}$；甲烷为限量反应物。
\textbf{8.} \ce{CH4} 0，\ce{H2O} \qty{2.00}{mol}，\ce{CO} \qty{1.00}{mol}，
\ce{H2} \qty{3.00}{mol}。
\textbf{9.} \qty{6.00}{mol}，对比投料的 \qty{4.00}{mol}：
该反应由两个分子生成四个分子。
\textbf{10.} \ce{H2O} 0.333，\ce{CO} 0.167，\ce{H2} 0.500。
\textbf{11.} 为保证甲烷全部反应，并为变换反应
留出水蒸气，其反应商因此降低。
\textbf{12.} \ce{CO} $1 - \xi$，\ce{H2O} $2 - \xi$，\ce{CO2} $\xi$，\ce{H2}
$3 + \xi$，任一反应进度下总量均为 6.00。
\textbf{13.} $Q = \dfrac{(p_{\ce{CO2}}/p^\circ)(p_{\ce{H2}}/p^\circ)}
{(p_{\ce{CO}}/p^\circ)(p_{\ce{H2O}}/p^\circ)}$，代入 $p_i = (n_i/6)p$：
其中因子 $p/(6p^\circ)$ 在分子、分母间相消（两侧各有两个气体
分子），余下所给出的表达式。
\textbf{14.} $\xi = 0$：$Q = 0 < 4.0$，正向进行。
\textbf{15.} $\xi(3 + \xi) = 4(1 - \xi)(2 - \xi)$，即 $3\xi^2 - 15\xi +
8 = 0$：$\xi = (15 - \sqrt{129})/6 = \qty{0.607}{mol}$；另一根
4.39，超过 $\xi_{\max} = 1$。
\textbf{16.} \ce{CO} \qty{0.393}{mol}，\ce{H2O} \qty{1.393}{mol}，
\ce{CO2} \qty{0.607}{mol}，\ce{H2} \qty{3.607}{mol}。
\textbf{17.} $\tau = 0.607/1.00 = 0.607$。
\textbf{18.} $\xi(3 + \xi) = 4(1 - \xi)(5 - \xi)$，$3\xi^2 - 27\xi + 20 =
0$，$\xi = \qty{0.814}{mol}$：水蒸气越多，转化率越高。
\textbf{19.} $K^\circ = (0.99 \times 3.99)/(0.01 \times 1.01) = 391$。
\textbf{20.} $Q(\xi)$ 单调递增，因此较大的 $K^\circ$ 在较大的
反应进度处达到。该反应的平衡常数在较低温度下较大，
温度下反应较慢：第一反应器在高温下运行、速度快，
第二反应器温度较低，用于完成转化（
$K^\circ$ 随温度的变化在本科二年级分册中讨论）。
\textbf{21.} $3.607/(6.00 - 1.393) = 3.607/4.607 = 0.783$。
\textbf{22.} \qty{0.393}{mol} 的 \ce{CO} 与 \qty{0.607}{mol} 的 \ce{CO2}。
\textbf{23.} $x(\ce{H2}) = 3.607/6.00 = \textbf{0.60}$：离开变换反应器的
十个分子中有六个为氢气。
\end{solution}
